% Нэр томьёоны тайлбар — засварлагдсан, бүлэглэгдсэн
% Encoding: UTF-8
% Compiler: XeLaTeX эсвэл LuaLaTeX
% Шаардлагатай package: longtable, booktabs, array

{\small
\renewcommand{\arraystretch}{1.2}
\begin{longtable}{p{0.24\textwidth} p{0.22\textwidth} p{0.10\textwidth} p{0.34\textwidth}}
\caption{Тайланд хэрэглэсэн гол нэр томьёоны тайлбар}\label{tab:glossary}\\
\toprule
\textbf{Англи нэр томьёо} & \textbf{Монгол орчуулга} & \textbf{Товчлол} & \textbf{Тайлбар} \\
\midrule
\endfirsthead
\toprule
\textbf{Англи нэр томьёо} & \textbf{Монгол орчуулга} & \textbf{Товчлол} & \textbf{Тайлбар} \\
\midrule
\endhead

% ============================================================
% А. ITS БА СУРГАЛТЫН ОНОЛ
% ============================================================
\multicolumn{4}{l}{\textbf{\textit{А. Хиймэл оюунт зааварлагч систем ба сургалтын онол}}} \\
\midrule

Computer-Assisted Instruction & Компьютерын тусламжтай сургалт & CAI & Компьютер ашиглан сургалтын агуулга хүргэх эрт үеийн арга зүй. \\

Computer-Based Learning System & Компьютерт суурилсан сургалтын систем & - & Компьютероор дамжуулан агуулга хүргэдэг уламжлалт сургалтын систем. \\

Intelligent Tutoring System & Хиймэл оюунт зааварлагч систем & ITS & Суралцагчийн мэдлэг, ахиц, алдаанд үндэслэн хувь хүнд тохирсон тайлбар, дасгал, зөвлөмж өгдөг сургалтын систем. \\

AI-powered ITS & Хиймэл оюунд суурилсан зааварлагч систем & - & LLM, RAG зэрэг орчин үеийн AI технологи ашигласан уян хатан зааварлагч систем. \\

Domain Model & Салбарын мэдлэгийн загвар & - & Юуг заахыг тодорхойлдог мэдлэгийн агуулга, ойлголт, дүрэм, даалгаврын бүтэц. \\

Student Model & Суралцагчийн загвар & - & ITS-ийн сонгодог нэршлээр суралцагчийн мэдлэгийн төлөвийг илэрхийлдэг бүрэлдэхүүн. \\

Learner Model & Суралцагчийн загвар & - & Суралцагчийн түвшин, хэрэгцээ, ахицад тулгуурласан динамик загвар. \\

Pedagogical Model & Сургалтын арга зүйн загвар & - & Яаж заах, ямар тайлбар ба дэмжлэг өгөхийг шийддэг загвар. \\

Tutoring Model & Зааварчилгааны загвар & - & Суралцагчид өгөх заавар, тусламж, зөвлөмжийн стратегийг тодорхойлдог хэсэг. \\

Interface Model & Харилцах орчны загвар & - & Суралцагч системтэй хэрхэн харилцахыг тодорхойлдог интерфейсийн түвшин. \\

User Interface & Хэрэглэгчийн интерфейс & UI & Хэрэглэгчтэй шууд харилцах дэлгэц, удирдлагын орчин. \\

Learner Profile & Суралцагчийн профайл & - & Суралцагчийн mastery, ахиц, сул ба хүчтэй тал, суралцах зан төлөвийг нэгтгэн харуулсан төлөвийн товч зураглал. \\

Open Learner Model & Нээлттэй суралцагчийн загвар & - & Суралцагчийн загварын төлөв, ахиц, сул ба хүчтэй талыг суралцагчид ил тод харуулдаг хандлага. \\

Black Box & Хаалттай загвар & - & Дотоод шийдвэр гаргалтын логик нь хэрэглэгчид ил тод биш, яагаад ийм үр дүн гарсныг тайлбарлахад хүндрэлтэй систем. \\

\midrule
% ============================================================
% Б. МЭДЛЭГИЙН АНГИЛАЛ (TAXONOMY)
% ============================================================
\multicolumn{4}{l}{\textbf{\textit{Б. Мэдлэгийн ангилал}}} \\
\midrule

Knowledge Taxonomy & Мэдлэгийн ангиллын тогтолцоо & - & Мэдлэгийг шаталсан бүтцээр зохион байгуулсан ангилал. \\

Body of Knowledge & Мэдлэгийн цогц & BOK & Тухайн мэргэжлийн салбарт багтах стандартчилагдсан мэдлэгийн хүрээ. \\

Domain & Домэйн & - & Системийн хамрах хамгийн дээд түвшний мэдлэгийн салбар. \\

Knowledge Area & Мэдлэгийн бүлэг & KA & Домэйн доторх томоохон агуулгын хэсэг. \\

Knowledge Unit & Мэдлэгийн нэгж & KU & Knowledge Area доторх нарийвчилсан агуулгын нэгж. \\

Topic & Сэдэв & - & Knowledge Unit доторх илүү нарийн агуулгын дэд хэсэг. \\

Learning Outcome & Суралцахуйн үр дүн & - & Суралцагч юуг мэдэж, хийж чаддаг болохыг тодорхойлсон үр дүн. \\

\midrule
% ============================================================
% В. СУРАЛЦАХУЙН ОНОЛ (MASTERY, ADAPTIVE)
% ============================================================
\multicolumn{4}{l}{\textbf{\textit{В. Суралцахуйн онол}}} \\
\midrule

Mastery Learning & Эзэмшилтийн сургалт & - & Агуулгыг хангалттай эзэмшсэний дараа дараагийн шат руу шилжүүлэх сургалтын хандлага. \\

Mastery Score & Эзэмшилтийн оноо & - & Суралцагч тухайн нэгжийг хэр зэрэг эзэмшсэнийг илэрхийлэх тоон үзүүлэлт. \\

Mastery Tracking & Эзэмшилтийн түвшний хяналт & - & Суралцагчийн KA/KU түвшний ахиц, эзэмшлийг оноогоор хянах механизм. \\

Closed Feedback Loop & Хаалттай буцаан холбооны давталт & - & Үнэлгээ, буцаан холбоо, засварлагч дэмжлэг, дахин үнэлгээг тасралтгүй давтан хэрэгжүүлдэг сургалтын мөчлөг. \\

Corrective Instruction & Засварлагч зааварчилгаа & - & Суралцагч шаардлагатай түвшинд хүрээгүй үед нэмэлт тайлбар, өөр жишээ, давтлагаар алдааг засах зорилготой сургалтын дэмжлэг. \\

Enrichment & Баяжуулсан сургалтын үйл ажиллагаа & - & Тухайн нэгжийг эзэмшсэн суралцагчдад илүү гүнзгийрүүлсэн, өргөтгөсөн агуулга ба даалгавар өгөх дэмжлэг. \\

Adaptive Learning & Дасан зохицох сургалт & - & Суралцагчийн түвшин, хэрэгцээнд тааруулан агуулга ба дэмжлэгийг өөрчилдөг сургалт. \\

Adaptive Decision & Дасан зохицох шийдвэр & - & Суралцагчийн төлөв, нотолгоо, суралцах зорилгод тулгуурлан дараагийн тайлбар, дасгал, зөвлөмжийг сонгох шийдвэрийн механизм. \\

Adaptive Recommendation & Дасан зохицох зөвлөмж & - & Суралцагчийн төлөвт тохируулан дараагийн агуулга, давтлага, дасгалыг санал болгох механизм. \\

Personalization & Хувь хүнд тохируулалт & - & Суралцагч бүрийн өгөгдөлд тулгуурлан агуулга, зөвлөмж, суралцах туршлагыг хувьчлан тааруулах үйл явц. \\

Personalized Learning & Хувь хүнд тохирсон сургалт & - & Суралцагчийн онцлог, ахиц, хэрэгцээнд тулгуурласан сургалтын хандлага. \\

Evidence & Нотолгоо & - & Дүгнэлт эсвэл хариултыг баталгаажуулахад ашиглах баримт, эх сурвалж, нотлох мэдээлэл. \\

Belief & Магадлалын итгэл & - & Суралцагч тухайн мэдлэгийг эзэмшсэн байх магадлалын талаарх системийн итгэлийн түвшин. \\

Learning Analytics & Суралцахуйн шинжилгээ & - & Суралцагчийн үйлдэл, ахиц, гүйцэтгэлийг өгөгдлөөр шинжлэх арга. \\

\midrule
% ============================================================
% Г. RAG БА LLM
% ============================================================
\multicolumn{4}{l}{\textbf{\textit{Г. RAG ба LLM}}} \\
\midrule

Large Language Model & Том хэлний загвар & LLM & Их хэмжээний өгөгдөл дээр сурсан, ойлгох ба үүсгэх чадвартай хэлний загвар. \\

Hallucination & Зохиомол мэдээллийн алдаа & - & LLM бодит эх сурвалжгүй мэдээллийг үнэмшилтэйгээр зохиох алдаа. \\

Retrieval-Augmented Generation & Хайлтаар баяжуулсан үүсгэлт & RAG & Гадаад эх сурвалжаас мэдээлэл татаж, түүнд тулгуурлан хариулт үүсгэдэг арга. \\

Retrieval-first Architecture & Эхлээд хайх архитектур & - & LLM-ийн дотоод мэдлэгээс өмнө эх сурвалж хайж контекст бүрдүүлдэг зарчим. \\

Retrieval & Холбогдох мэдээлэл татах & - & Хэрэглэгчийн асуултад холбоотой эх сурвалжийн хэсгийг олж татах үйл явц. \\

Augmentation & Контекст нөхөн бүрдүүлэлт & - & Олдсон эх сурвалжийг асуулттай нэгтгэн хариултын контекст болгох алхам. \\

Generation & Хариулт үүсгэлт & - & LLM өгөгдсөн контекст дээр тулгуурлан тайлбар, хариулт гаргах үйл явц. \\

Citation & Эшлэл & - & Хариултад ашигласан эх сурвалжийн холбоос, эшлэл, хуудасны мэдээлэл. \\

Source-grounded Answer & Эх сурвалжид суурилсан хариулт & - & Хэрэглэгчийн оруулсан эсвэл баталгаатай эх сурвалж дээр суурилсан хариулт. \\

\midrule
% ============================================================
% Д. БАРИМТ БОЛОВСРУУЛАЛТ (INGESTION PIPELINE)
% ============================================================
\multicolumn{4}{l}{\textbf{\textit{Д. Баримт боловсруулалт}}} \\
\midrule

Document Ingestion Pipeline & Баримт боловсруулах дамжуулга & - & Баримтыг хүлээн авах, задлах, хэсэглэх, вектор болгох, индексжүүлэх урсгал. \\

PDF Parsing & PDF задлалт & - & PDF файлаас текст, хуудасны мэдээлэл гарган авах үйл явц. \\

Chunk & Текстийн хэсэг & - & Урт баримтаас хайлт ба боловсруулалтад зориулан хуваасан жижиг нэгж. \\

Chunking & Текст хэсэглэх & - & Текстийг жижиг, утгын хувьд удирдах боломжтой хэсгүүдэд хуваах арга. \\

Recursive Character Text Splitter & Рекурсив тэмдэгт хуваагч & - & Текстийг шаталсан байдлаар (параграф $\rightarrow$ өгүүлбэр $\rightarrow$ тэмдэгт) хуваадаг түгээмэл chunking хэрэгсэл. \\

Embedding & Вектор дүрслэл & - & Текстийн утга санааг олон хэмжээст тоон вектор хэлбэрт хувиргах процесс. \\

Embedding Model & Вектор дүрслэл үүсгэх загвар & - & Текстээс вектор дүрслэл гаргадаг загвар. \\

Vector Database & Вектор өгөгдлийн сан & - & Вектор хэлбэрийн өгөгдөл хадгалж, утгын ойролцоо хайлт хийдэг сан. \\

Vector Search & Вектор хайлт & - & Векторуудын ижил төстэй байдлаар хайлт хийх арга. \\

Semantic Search & Утга санааны хайлт & - & Түлхүүр үг бус утгын ойролцоо байдлаар мэдээлэл олох хайлт. \\

Top-k & Шилдэг k үр дүн & - & Хайлтаар хамгийн ойролцоо гэж сонгосон k ширхэг үр дүн. \\

Reranking & Дахин эрэмбэлэлт & - & Эхний хайлтын үр дүнг илүү оновчтой дарааллаар дахин эрэмбэлэх арга. \\

OCR & Оптик тэмдэгт таних & OCR & Зураг эсвэл сканердсан баримтаас текст таньж гаргах технологи. \\

Batch Processing & Багц боловсруулалт & - & Нэг удаад олон нэгж өгөгдлийг хамтад нь боловсруулах арга. \\

Streaming Processing & Урсгал боловсруулалт & - & Өгөгдлийг тасралтгүй урсгалаар бага багаар боловсруулах арга. \\

Metadata & Мета өгөгдөл & - & Баримтын нэр, төрөл, хэмжээ, хуудас зэргийг тайлбарлах өгөгдөл. \\

Source Type & Эх сурвалжийн төрөл & - & Баримтын төрлийг илэрхийлэх талбар; жишээ нь PDF, plain text. \\

\midrule
% ============================================================
% Е. ТЕХНОЛОГИЙН СТЕК
% ============================================================
\multicolumn{4}{l}{\textbf{\textit{Е. Технологийн стек}}} \\
\midrule

Technology Stack & Технологийн стек & - & Системийг хэрэгжүүлэхэд ашиглаж буй технологийн цогц багц. \\

Frontend & Урд талын хөгжүүлэлт & - & Хэрэглэгчид харагдах веб интерфейсийн давхарга. \\

Backend & Арын талын хөгжүүлэлт & - & Сервер талын логик, өгөгдөл боловсруулалт, API удирддаг давхарга. \\

Application Programming Interface & Хэрэглээний программчлалын интерфейс & API & Программ хооронд өгөгдөл, үйлчилгээ солилцох дүрэм, интерфейс. \\

AI Orchestration & AI урсгалын зохион байгуулалт & - & Олон алхамт AI pipeline-ийн дараалал, төлөв, уялдааг удирдах арга. \\

AI Orchestration & AI урсгалын зохион байгуулалт & - & Олон алхамт AI pipeline-ийн дараалал, төлөв, уялдааг удирдах арга. \\

Weaviate & Weaviate & - & Вектор хайлт, retrieval-д ашигласан vector database хөдөлгүүр. \\

BGE-M3 & BGE-M3 & - & Олон хэл дэмждэг, dense/sparse/multi-vector хайлтыг нэгтгэсэн embedding загвар. \\

PostgreSQL & PostgreSQL & - & Харилцаат өгөгдлийн сангийн удирдлагын систем. \\

Prisma & Prisma & ORM & TypeScript орчинд өгөгдлийн сантай ажиллах ORM хэрэгсэл. \\

Object-Relational Mapping & Объект-харилцаат зураглал & ORM & Программын объект ба өгөгдлийн сангийн хүснэгтийг холбодог арга. \\

Next.js & Next.js & - & React дээр суурилсан full-stack веб framework. \\

React & React & - & Component-д суурилсан веб интерфейс хөгжүүлэх JavaScript сан. \\

TypeScript & TypeScript & TS & Төрөлтэй JavaScript өргөтгөл хэл. \\

Node.js & Node.js & - & JavaScript/TypeScript кодыг сервер талд ажиллуулах runtime орчин. \\

Worker & Туслах боловсруулагч процесс & - & Хүнд даацын баримт боловсруулалтыг үндсэн процессоос тусад нь гүйцэтгэх процесс. \\

BullMQ & BullMQ & - & Node.js дээрх job queue сан; даалгаврыг дараалалд оруулж боловсруулна. \\

Redis & Redis & - & Хурдан in-memory key-value store; queue ба cache-д өргөн хэрэглэгддэг. \\

\midrule
% ============================================================
% Ж. СИСТЕМИЙН ШААРДЛАГА БА ДИЗАЙН
% ============================================================
\multicolumn{4}{l}{\textbf{\textit{Ж. Системийн шаардлага ба дизайн}}} \\
\midrule

Business Requirements & Бизнесийн шаардлага & - & Систем бизнесийн түвшинд ямар зорилго, үнэ цэнэ хангахыг тодорхойлсон шаардлага. \\

User Requirements & Хэрэглэгчийн шаардлага & - & Хэрэглэгч системээс юу хүсэж байгааг тодорхойлсон шаардлага. \\

Functional Requirements & Функциональ шаардлага & - & Систем ямар үйлдэл, боломж хэрэгжүүлэхийг тодорхойлсон шаардлага. \\

Non-Functional Requirements & Функциональ бус шаардлага & - & Гүйцэтгэл, найдвартай байдал, өргөтгөх боломж зэрэг чанарын шаардлага. \\

Use Case & Хэрэглээний тохиолдол & - & Хэрэглэгч системтэй хэрхэн харилцахыг дүрсэлсэн загвар. \\

ER Diagram & Entity-Relationship диаграмм & ERD & Өгөгдлийн сангийн entity ба холбоосыг дүрслэх диаграмм. \\

CRUD & Үүсгэх, унших, шинэчлэх, устгах & CRUD & Өгөгдөлтэй ажиллах үндсэн дөрвөн үйлдэл. \\

Minimum Viable Product & Хамгийн бага хэрэгжихүйц бүтээгдэхүүн & MVP & Гол үнэ цэнийг шалгахад шаардлагатай хамгийн бага боломжуудтай хувилбар. \\

Dashboard & Хяналтын самбар & - & Ахиц, оноо, үзүүлэлтийг нэг дор харуулах хэрэглэгчийн дэлгэц. \\

\midrule
% ============================================================
% З. БУСАД
% ============================================================
\multicolumn{4}{l}{\textbf{\textit{З. Бусад}}} \\
\midrule

Plain Text & Энгийн текст & - & Тэмдэглэгээ, форматгүй цэвэр текстэн өгөгдөл. \\

Multimodal Input & Олон хэлбэрийн оролт & - & Текст, дуу, дүрс зэрэг олон эх үүсвэрийн оролт. \\

Speech-to-Text & Яриаг текст болгох & STT & Ярианы дохиог бичвэр болгон хувиргах технологи. \\

Design Science & Дизайн шинжлэх ухаан & - & Бодит асуудлыг шийдэх артефакт бүтээж, түүнийг үнэлэх замаар мэдлэг бий болгодог судалгааны хандлага. \\

\bottomrule
\end{longtable}
}